\documentclass{amsart}
% For dealing with references we use the comment environment
\usepackage{verbatim}
\newenvironment{reference}{\comment}{\endcomment}
%\newenvironment{reference}{}{}
\newenvironment{slogan}{\comment}{\endcomment}
\newenvironment{history}{\comment}{\endcomment}

% For commutative diagrams we use Xy-pic
\usepackage[all]{xy}

% We use 2cell for 2-commutative diagrams.
\xyoption{2cell}
\UseAllTwocells

% We use multicol for the list of chapters between chapters
\usepackage{multicol}

% This is generally recommended for better output
\usepackage{lmodern}
\usepackage[T1]{fontenc}

% For cross-file-references

% Package for hypertext links:
\usepackage{hyperref}

% For any local file, say "hello.tex" you want to link to please

% Theorem environments.
%
\theoremstyle{plain}
\newtheorem{theorem}[subsection]{Theorem}
\newtheorem{proposition}[subsection]{Proposition}
\newtheorem{lemma}[subsection]{Lemma}

\theoremstyle{definition}
\newtheorem{definition}[subsection]{Definition}
\newtheorem{example}[subsection]{Example}
\newtheorem{exercise}[subsection]{Exercise}
\newtheorem{situation}[subsection]{Situation}

\theoremstyle{remark}
\newtheorem{remark}[subsection]{Remark}
\newtheorem{remarks}[subsection]{Remarks}

\numberwithin{equation}{subsection}

% Macros
%
\def\lim{\mathop{\mathrm{lim}}\nolimits}
\def\colim{\mathop{\mathrm{colim}}\nolimits}
\def\Spec{\mathop{\mathrm{Spec}}}
\def\Hom{\mathop{\mathrm{Hom}}\nolimits}
\def\Ext{\mathop{\mathrm{Ext}}\nolimits}
\def\SheafHom{\mathop{\mathcal{H}\!\mathit{om}}\nolimits}
\def\SheafExt{\mathop{\mathcal{E}\!\mathit{xt}}\nolimits}
\def\Sch{\mathit{Sch}}
\def\Mor{\mathop{\mathrm{Mor}}\nolimits}
\def\Ob{\mathop{\mathrm{Ob}}\nolimits}
\def\Sh{\mathop{\mathit{Sh}}\nolimits}
\def\NL{\mathop{N\!L}\nolimits}
\def\CH{\mathop{\mathrm{CH}}\nolimits}
\def\proetale{{pro\text{-}\acute{e}tale}}
\def\etale{{\acute{e}tale}}
\def\QCoh{\mathit{QCoh}}
\def\Ker{\mathop{\mathrm{Ker}}}
\def\Im{\mathop{\mathrm{Im}}}
\def\Coker{\mathop{\mathrm{Coker}}}
\def\Coim{\mathop{\mathrm{Coim}}}

% Boxtimes
%
\DeclareMathSymbol{\boxtimes}{\mathbin}{AMSa}{"02}

%
% Macros for moduli stacks/spaces
%
\def\QCohstack{\mathcal{QC}\!\mathit{oh}}
\def\Cohstack{\mathcal{C}\!\mathit{oh}}
\def\Spacesstack{\mathcal{S}\!\mathit{paces}}
\def\Quotfunctor{\mathrm{Quot}}
\def\Hilbfunctor{\mathrm{Hilb}}
\def\Curvesstack{\mathcal{C}\!\mathit{urves}}
\def\Polarizedstack{\mathcal{P}\!\mathit{olarized}}
\def\Complexesstack{\mathcal{C}\!\mathit{omplexes}}
% \Pic is the operator that assigns to X its picard group, usage \Pic(X)
% \Picardstack_{X/B} denotes the Picard stack of X over B
% \Picardfunctor_{X/B} denotes the Picard functor of X over B
\def\Pic{\mathop{\mathrm{Pic}}\nolimits}
\def\Picardstack{\mathcal{P}\!\mathit{ic}}
\def\Picardfunctor{\mathrm{Pic}}
\def\Deformationcategory{\mathcal{D}\!\mathit{ef}}

\setlength{\emergencystretch}{3em}
\begin{document}
\title[Stacks Project, Tag 03I1]{479 --- Every quasi-finite morphism becomes finite over a dense open of its target}
\maketitle
\noindent Copyright (C) 2005--2025 Johan de Jong. Stacks Project authors. Source: \href{https://stacks.math.columbia.edu/tag/03I1}{Tag 03I1}.

\noindent Editorial difficulty note (not author text). Difficulty H2: The morphism need not be separated, so generic finiteness of affine charts is insufficient. The proof removes images of overlap boundaries to make the chart intersections closed and recover separatedness, then descends properness from a finite chart cover to establish finiteness..

\noindent Editorial provenance note (not author text). History: excerpt from revision \texttt{a0444\allowbreak 6e57e\allowbreak c1fbc\allowbreak 252a8\allowbreak 71afc\allowbreak ec775\allowbreak 2fb28\allowbreak 07b14}, more-morphisms.tex, lines 12677--12758. Reference presentation was adapted; original-author context and core support blocks were retained or added. The mathematical wording is unchanged. Licensed under GNU FDL 1.2 or later; no invariant sections or cover texts. See ../licenses/COPYING. Context blocks reproduce author definitions and supporting results; they are not additional counted problems.

\begin{lemma}
\label{lemma-quasi-finite-finite-over-dense-open}
Let $f : Y \to X$ be a quasi-finite morphism.
There exists a dense open $U \subset X$ such that
$f|_{f^{-1}(U)} : f^{-1}(U) \to U$ is finite.
\end{lemma}

\begin{proof}
If $U_i \subset X$, $i \in I$ is a collection of opens such that the
restrictions $f|_{f^{-1}(U_i)} : f^{-1}(U_i) \to U_i$ are finite,
then with $U = \bigcup U_i$ the restriction $f|_{f^{-1}(U)} : f^{-1}(U) \to U$
is finite, see
Morphisms, Lemma \href{https://stacks.math.columbia.edu/tag/01WI}{lemma finite local (Tag 01WI)}.
Thus the problem is local on $X$ and we may assume that $X$ is affine.

\medskip\noindent
Assume $X$ is affine.
Write $Y = \bigcup_{j = 1, \ldots, m} V_j$ with $V_j$ affine.
This is possible since $f$ is quasi-finite and hence
in particular quasi-compact. Each $V_j \to X$ is quasi-finite
and separated. Let $\eta \in X$ be a generic point of an irreducible
component of $X$. We see from
Morphisms, Lemmas
\href{https://stacks.math.columbia.edu/tag/02NH}{lemma quasi finite (Tag 02NH)} and \href{https://stacks.math.columbia.edu/tag/02NW}{lemma generically finite (Tag 02NW)}
that there exists an open neighbourhood $\eta \in U_\eta$ such that
$f^{-1}(U_\eta) \cap V_j \to U_\eta$ is finite. We may choose $U_\eta$ such
that it works for each $j = 1, \ldots, m$.
Note that the collection of generic points of $X$ is dense in $X$.
Thus we see there exists a dense open $W = \bigcup_\eta U_\eta$
such that each $f^{-1}(W) \cap V_j \to W$ is finite.
It suffices to show that there exists a dense open $U \subset W$
such that $f|_{f^{-1}(U)} : f^{-1}(U) \to U$ is finite.
Thus we may replace $X$ by an affine open subscheme of $W$ and
assume that each $V_j \to X$ is finite.

\medskip\noindent
Assume $X$ is affine, $Y = \bigcup_{j = 1, \ldots, m} V_j$ with $V_j$ affine,
and the restrictions $f|_{V_j} : V_j \to X$ are finite.
Set
$$
\Delta_{ij} =
\Big(\overline{V_i \cap V_j} \setminus V_i \cap V_j\Big) \cap V_j.
$$
This is a nowhere dense closed subset of $V_j$ because it is the boundary
of the open subset $V_i \cap V_j$ in $V_j$. By
Morphisms, Lemma \href{https://stacks.math.columbia.edu/tag/03HX}{lemma image nowhere dense finite (Tag 03HX)}
the image $f(\Delta_{ij})$ is a nowhere dense closed subset of $X$. By
Topology, Lemma \href{https://stacks.math.columbia.edu/tag/03HO}{lemma nowhere dense (Tag 03HO)}
the union $T = \bigcup f(\Delta_{ij})$ is a nowhere dense closed
subset of $X$. Thus $U = X \setminus T$ is a dense open subset of $X$.
We claim that $f|_{f^{-1}(U)} : f^{-1}(U) \to U$ is finite.
To see this let $U' \subset U$ be an affine open.
Set $Y' = f^{-1}(U') = U' \times_X Y$,
$V_j' = Y' \cap V_j = U' \times_X V_j$. Consider the restriction
$$
f' = f|_{Y'} : Y' \longrightarrow U'
$$
of $f$. This morphism now has the property that
$Y' = \bigcup_{j = 1, \ldots, m} V'_j$ is an affine open covering,
each $V'_j \to U'$ is finite, and $V_i' \cap V_j'$ is (open and) closed
both in $V'_i$ and $V'_j$. Hence $V_i' \cap V_j'$ is affine, and the map
$$
\mathcal{O}(V'_i) \otimes_{\mathbf{Z}} \mathcal{O}(V'_j)
\longrightarrow
\mathcal{O}(V'_i \cap V'_j)
$$
is surjective. This implies that $Y'$ is separated, see
Schemes, Lemma \href{https://stacks.math.columbia.edu/tag/01KP}{lemma characterize separated (Tag 01KP)}.
Finally, consider the commutative diagram
$$
\xymatrix{
\coprod_{j = 1, \ldots, m} V'_j \ar[rd] \ar[rr] & & Y' \ar[ld] \\
& U' &
}
$$
The south-east arrow is finite, hence proper, the horizontal arrow is
surjective, and the south-west arrow is separated. Hence by
Morphisms, Lemma \href{https://stacks.math.columbia.edu/tag/03GN}{lemma image proper is proper (Tag 03GN)}
we conclude that $Y' \to U'$ is proper. Since it is also quasi-finite,
we see that it is finite by Lemma \href{https://stacks.math.columbia.edu/tag/02LS}{lemma characterize finite (Tag 02LS)},
and we win.
\end{proof}
\end{document}
